
Paraphrasing defeats contamination detection---but contamination does leave behavioral signatures. Through controlled experiments, this work demonstrates that contamination creates training exposure (31.1\% accuracy gain), paraphrase-augmented training produces representation invariance (MPS difference 0.065, d=0.52), and invariance manifests as confidence uniformity (r=$-0.517)$. The causal chain from contamination to confidence behavior is validated.

However, SSI = 1/variance fails as a practical detector (AUC=0.506 on real data). The inverse-variance formulation amplifies noise rather than signal. The contributions are: (1) mechanism validation demonstrating that contamination produces detectable behavioral changes, (2) documentation of inverse-variance metric failure guiding future work toward alternative formulations, and (3) identification of the simulation-reality gap highlighting the necessity of real-data validation.

Alternative metrics leveraging the validated mechanism---entropy-based, coefficient of variation, or rank-based---offer paths forward for paraphrase-resistant contamination detection.
